\PDEIndexMark{idx:analytic-fredholm-main-theorem-introduction}现在证明一个有时称为\emph{解析 Fredholm 定理}的结果. 它是两个重要结果的基础: Fredholm 二择一定理 (Theorem~\ref{thm:fredholm-alternative-compact}) 和 Hilbert-Schmidt 定理 (Theorem~\ref{thm:hilbert-schmidt-spectral}).

\begin{Theorem}{解析 Fredholm 定理}
\label{thm:analytic-fredholm}
设 $G\subset\mathbb{C}$ 为区域. 假设映射
\begin{equation}\label{eq:analytic-compact-fredholm-operator-family}
\mathbb{C}\supset G\ni\lambda\mapsto B(\lambda)\in\mathcal{L}(H)
\end{equation}
在 $G$ 上解析, 且对每个 $\lambda\in G$, $B(\lambda)$ 都是紧算子. 那么以下两者之一成立:\begin{enumerate}\item 对任何 $\lambda\in G$, $(I-B(\lambda))^{-1}$ 均不存在, 或者\item 对每个 $\lambda\in G\backslash S$, $(I-B(\lambda))^{-1}$ 均存在, 其中 $S$ 是 $G$ 中的离散集 (即在 $G$ 中没有极限点). 此时函数 $\lambda\mapsto(I-B(\lambda))^{-1}$ 在 $G\backslash S$ 上解析, 且若 $\lambda\in S$, 则 $B(\lambda)\psi=\psi$ 有有限维解族.\end{enumerate}
\end{Theorem}

\begingroup
\postdisplaypenalty=10000
\renewcommand{\endProof}{\par\nobreak\hfill$\blacksquare$\par\vspace{0.1cm}}
\begin{Proof}
只在某点 $\lambda_0\in G$ 的邻域内给出证明. 可以用标准的连通性论证将结果推广到整个 $G$.

给定 $\lambda_0\in G$. 由于 $\lambda\mapsto B(\lambda)$ 连续, 可以选择 $r>0$, 使得对圆盘 $D_r=\{\lambda\in G\mid|\lambda-\lambda_0|<r\}$ 中的所有 $\lambda$, 均有
\begin{equation}\label{eq:analytic-fredholm-local-continuity-bound}
\|B(\lambda)-B(\lambda_0)\|<\dfrac{1}{2}
\end{equation}
利用 Theorem~\ref{thm:compact-hilbert-finite-rank-approximation} 中的构造, 可知存在有限秩算子 $B_N$, 使得
\begin{equation}\label{eq:analytic-fredholm-finite-rank-approximation-bound}
\|B_N-B(\lambda_0)\|<1/2.
\end{equation}
现在利用 Theorem~\ref{thm:bounded-operator-analytic-resolvent} 的证明中的几何级数方法, 请读者验证
\begin{equation}\label{eq:analytic-fredholm-small-remainder-inverse}
(I-B(\lambda)+B_N)^{-1}
\end{equation}
作为有界算子存在, 且在 $D_r$ 上解析.

现在令
\begin{equation}\label{eq:analytic-fredholm-finite-rank-reduction}
F(\lambda):=B_N\circ(I-B(\lambda)+B_N)^{-1}.
\end{equation}
注意
\begin{equation}\label{eq:analytic-fredholm-operator-factorization}
(I-B(\lambda))=(I-F(\lambda))(I-B(\lambda)+B_N).
\end{equation}
因而 $I-B(\lambda)$ 可逆, 当且仅当 $I-F(\lambda)$ 可逆. 然而, $F$ 具有有限秩, 因而由 Lemma~\ref{lem:analytic-finite-rank-alternative}, 要么对任何 $\lambda\in G$, $I-F(\lambda)$ 都不可逆, 要么它在离散集 $S\subset G$ 之外可逆.

$B(\lambda)\psi=\psi$ 的解空间有限维这一点, 可由 $B(\lambda)$ 的紧性得到, 证明留给读者 (Problem~\ref{ex:compact-nonzero-eigenspace-finite-dimensional}). 证明完毕.
\end{Proof}
\endgroup

现在利用解析 Fredholm 定理, 导出以下关于 Hilbert 空间上紧算子的谱的刻画.

\begin{Theorem}{Fredholm 二择一定理}
\label{thm:fredholm-alternative-compact}
设 $A\in\mathcal{L}(H)$ 为紧算子. 那么 $\sigma(A)$ 为紧集, 除可能的 $\lambda=0$ 外没有极限点. 此外, 给定任意 $\lambda\in\mathbb{C}\backslash\{0\}$, 以下两者之一成立:\begin{enumerate}\item $\lambda\in\rho(A)$, 或者\item $\lambda\in\sigma_p(A)$ 为有限重数的特征值.\end{enumerate}
\end{Theorem}

\begingroup
\postdisplaypenalty=10000
\renewcommand{\endProof}{\par\nobreak\hfill$\blacksquare$\par\vspace{0.1cm}}
\begin{Proof}
令 $G=\mathbb{C}\backslash\{0\}$, 并令
\begin{equation}\label{eq:compact-fredholm-scaled-operator-family}
B(\lambda)=\dfrac{1}{\lambda}A.
\end{equation}
然后注意到
\begin{equation}\label{eq:compact-fredholm-scaled-inverse-identity}
(\lambda I-A)^{-1}=\dfrac{1}{\lambda}\left(I-\dfrac{1}{\lambda}A\right)^{-1}=\dfrac{1}{\lambda}(I-B(\lambda))^{-1}.
\end{equation}
由 Theorem~\ref{lem:analytic-finite-rank-alternative} 直接得到结论.
\end{Proof}
\endgroup
